\chapter{\bilingual{Requisitos de usuario}{User Requirements}}

\section{\bilingual{Alcance del sistema}{System Scope}}

\bilingual{%
  Describe aquí el alcance del sistema a desarrollar: qué incluye y qué
  queda fuera del proyecto.
}{%
  Describe here the scope of the system to be developed: what is
  included and what is left out of the project.
}

\section{\bilingual{Requisitos de usuario}{User Requirements}}

\bilingual{%
  Los requisitos de usuario se numeran con el prefijo \textbf{RQ},
  pudiendo anidarse en subrequisitos:
}{%
  User requirements are numbered with the \textbf{RQ} prefix, and can be
  nested into sub-requirements:
}

\begin{requirements}
  \item \bilingual{Requisito 1}{Requirement 1}
  \begin{requirements}
    \item \bilingual{Subrequisito 1}{Sub-requirement 1}
    \item \bilingual{Subrequisito 2}{Sub-requirement 2}
    \begin{requirements}
      \item \bilingual{Subsubrequisito 1}{Sub-sub-requirement 1}
      \item \bilingual{Subsubrequisito 2}{Sub-sub-requirement 2}
    \end{requirements}
  \end{requirements}
  \item \bilingual{Requisito 2}{Requirement 2}
  \item \bilingual{Requisito 3}{Requirement 3}
  \item \bilingual{Requisito 4}{Requirement 4}
\end{requirements}

\section{\bilingual{Alternativas}{Alternatives}}

\bilingual{%
  Análisis de las alternativas consideradas para resolver el problema
  planteado, incluyendo metodologías ágiles como Scrum~\cite{scrumguide}.
}{%
  Analysis of the alternatives considered to solve the stated problem,
  including agile methodologies such as Scrum~\cite{scrumguide}.
}
